\SourceLeaf{66}%
tate deminutae divisibiles erunt per $1-x^p$; quamobrem in hoc casu etiam $f(x^n)-p$ per $1-x^p$ et proin etiam per $\frac{1-x^p}{1-x}$ divisibiles erit.

\Article{3}

\textsc{Theorema.}\quad \textit{Statuendo}

\[
\begin{aligned}
x-x^\alpha+x^{\alpha\alpha}-x^{\alpha^3}+x^{\alpha^4}-\text{etc.}-x^{\alpha^{p-2}}=\xi
\end{aligned}
\]

\textit{erit $\xi\xi\mp p$ divisibilis per $\frac{1-x^p}{1-x}$, accepto signo superiori, quoties $p$ est formae} \textit{$4k+1$, inferiori, quoties $p$ est formae $4k+3$.}

\textsc{Demonstr.}\quad Facile perspicietur, ex $p-1$ functionibus hisce

\[
\begin{aligned}
&+x\xi-xx+x^{\alpha+1}-x^{\alpha\alpha+1}\\&\qquad+\text{etc.}+x^{\alpha^{p-2}+1}\\
&-x^\alpha\xi-x^{2\alpha}+x^{\alpha\alpha+\alpha}-x^{\alpha^3+\alpha}\\&\qquad+\text{etc.}+x^{\alpha^{p-1}+\alpha}\\
&+x^{\alpha\alpha}\xi-x^{2\alpha\alpha}+x^{\alpha^3+\alpha\alpha}-x^{\alpha^4+\alpha\alpha}\\&\qquad+\text{etc.}+x^{\alpha^p+\alpha\alpha}\\
&-x^{\alpha^3}\xi-x^{2\alpha^3}+x^{\alpha^4+\alpha^3}-x^{\alpha^5+\alpha^3}\\&\qquad+\text{etc.}+x^{\alpha^{p+1}+\alpha^3}
\end{aligned}
\]

etc. usque ad

\[
\begin{aligned}
-x^{\alpha^{p-2}}\xi-x^{2\alpha^{p-2}}+x^{\alpha^{p-1}+\alpha^{p-2}}-x^{\alpha^p+\alpha^{p-2}}\\&\qquad+\text{etc.}+x^{\alpha^{2p-4}+\alpha^{p-2}}
\end{aligned}
\]

primam fieri $=0$, singulas reliquas autem per $1-x^p$ divisibiles. Quare per $1-x^p$ etiam divisibilis erit omnium summa, quae colligitur

\[
\begin{aligned}
&=\xi\xi-(f(xx)-1)+(f(x^{\alpha+1})-1)-(f(x^{\alpha\alpha+1})-1)\\&\qquad+(f(x^{\alpha^3+1})-1)-\text{etc.}\\
&\qquad+(f(x^{\alpha^{p-2}+1})-1)\\
&=\xi\xi-f(xx)+f(x^{\alpha+1})-f(x^{\alpha\alpha+1})\\&\qquad+f(x^{\alpha^3+1})-\text{etc.}+f(x^{\alpha^{p-2}+1})=\Omega
\end{aligned}
\]

Erit itaque haecce expressio $\Omega$ etiam divisibilis per $\frac{1-x^p}{1-x}$. Iam inter exponentes $2,\ \alpha+1,\ \alpha\alpha+1,\ \alpha^3+1\dotsfive\alpha^{p-2}+1$ unicus tantum erit divisibilis per $p$, puta $\alpha^{\frac12(p-1)}+1$, unde per art. praec. singulae partes expressionis $\Omega$ hae

\[
\begin{aligned}
f(xx),\quad f(x^{\alpha+1}),\quad f(x^{\alpha\alpha+1}),\quad (fx^{\alpha^3+1})\quad\text{etc.}
\end{aligned}
\]

excepto solo termino $f(x^{\alpha^{\frac12(p-1)}+1})$, divisibiles erunt per $\frac{1-x^p}{1-x}$. Istas itaque partes delere licebit, ita ut per $\frac{1-x^p}{1-x}$ etiam divisibilis maneat functio

\[
\begin{aligned}
\xi\xi\mp f(x^{\alpha^{\frac12(p-1)}+1})
\end{aligned}
\]

\SourceLeaf{67}%
ubi signum superius vel inferius valebit, prout $p$ est formae $4k+1$ vel formae $4k+3$. Et quum insuper $f(x^{\alpha^{\frac12(p-1)}+1})-p$ divisibilis sit per $\frac{1-x^p}{1-x}$, erit etiam $\xi\xi\mp p$ per $\frac{1-x^p}{1-x}$ divisibilis. Q. E. D.

Ne duplex signum ullam ambiguitatem adducere possit, per $\varepsilon$ numerum $+1$ vel $-1$ denotabimus, prout $p$ est formae $4k+1$ vel $4k+3$. Erit itaque $\frac{(1-x)(\xi\xi-\varepsilon p)}{1-x^p}$ functio integra ipsius $x$, quam per $Z$ designabimus.

\Article{4}

Sit $q$ numerus positivus impar, adeoque $\frac12(q-1)$ integer. Erit itaque $(\xi\xi)^{\frac12(q-1)}-(\varepsilon p)^{\frac12(q-1)}$ divisibilis per $\xi\xi-\varepsilon p$. et proin etiam per $\frac{1-x^p}{1-x}$. Statuamus $\varepsilon^{\frac12(q-1)}=\delta$, atque

\[
\begin{aligned}
\xi^{q-1}-\delta p^{\frac12(q-1)}=\frac{1-x^p}{1-x}\cdot Y
\end{aligned}
\]

eritque $Y$ functio integra ipsius $x$, atque $\delta=\mathord{.}+1$, quoties unus numerorum $p,q$, sive etiam uterque, est formae $4k+1$; contra erit $\delta=-1$, quoties uterque $p,q$ est formae $4k+3$.

\Article{5}

Iam supponamus, $q$ quoque esse numerum primum (a $p$ diversum) patetque per theorema in \textit{Disquisitionibus Arithmeticis} art. 51 demonstratum,

\[
\begin{aligned}
\mathord{.}\quad\xi^q-(x^q-x^{q\alpha}+x^{q\alpha\alpha}-x^{q\alpha^3}+\text{etc.}-x^{q\alpha^{p-2}})
\end{aligned}
\]

divisibilem fieri per $q$, sive formae $qX$, ita ut $X$ sit functio integra ipsius $x$ etiam respectu coëfficientium numericorum (quod etiam de functionibus reliquis integris hic occurrentibus $Z,Y,W$ subintelligendum est). Designemus pro modulo $p$ atque radice primitiva $\alpha$ indicem numeri $q$ per $\mu$, i. e. sit $q\equiv\alpha^\mu\ (\mathrm{mod.}\ p)$. Erunt itaque numeri $q,\ q\alpha,\ q\alpha\alpha,\ q\alpha^3\dotsfive q\alpha^{p-2}$ secundum modulum $p$ resp. congrui numeris $\alpha^\mu,\allowbreak\ \alpha^{\mu+1},\allowbreak\ \alpha^{\mu+2}\dotsfour\alpha^{p-2},\allowbreak\ 1,\allowbreak\ \alpha,\allowbreak\ \alpha\alpha\dotsfour\alpha^{\mu-1}$, adeoque

\[
\begin{aligned}
&x^q-x^{\alpha^\mu}\\
&x^{q\alpha}-x^{\alpha^{1+1}}\\
&x^{q\alpha\alpha}-x^{\alpha^{\mu+2}}\\
&x^{q\alpha^3}-x^{\alpha^{\mu+3}}\\
&\hspace{9mm}\begin{gathered}.\\.\\.\\.\end{gathered}
\end{aligned}
\]

\SourceLeaf{68}%
\[
\begin{aligned}
&x^{q\alpha^{p-\mu-2}}-x^{\alpha^{p-2}}\\
&x^{q\alpha^{p-\mu-1}}-x\\
&x^{q\alpha^{p-\mu}}-x^\alpha\\
&x^{q\alpha^{p-\mu+1}}-x^{\alpha\alpha}\\
&\hspace{9mm}\begin{gathered}.\\.\\.\\.\end{gathered}\\
&x^{q\alpha^{p-2}}-x^{\alpha^{\mu-1}}
\end{aligned}
\]

per $1-x^p$ divisibiles. Quibus quantitatibus, alternis vicibus positive et negative sumtis atque summatis, patet, per $1-x^p$ divisibilem esse functionem

\[
\begin{aligned}
x^q-x^{q\alpha}+x^{q\alpha\alpha}-x^{q\alpha^3}+\text{etc.}-x^{q\alpha^{p-2}}\mp\xi
\end{aligned}
\]

valente signo superiori vel inferiori, prout $\mu$ par sit vel impar, i. e. prout $q$ sit residuum quadraticum ipsius $p$ vel non-residuum. Statuemus itaque

\[
\begin{aligned}
x^q-x^{q\alpha}+x^{q\alpha\alpha}- x^{q\alpha^3}+\text{etc.}-x^{q\alpha^{p-2}}-\gamma\xi\\&=(1-x^p)W
\end{aligned}
\]

faciendo $\gamma=+1$, vel $\gamma=-1$, prout $q$ est residuum quadraticum ipsius $p$ vel non-residuum, patetque, $W$ fieri functionem integram.

\Article{6}

His ita praeparatis, e combinatione aequationum praecedentium deducimus

\[
\begin{aligned}
q\xi X=\varepsilon p(\delta p^{\frac12(q-1)}-\gamma)\\&\qquad+\frac{1-x^p}{1-x}\cdot(Z(\delta p^{\frac12(q-1)}-\gamma)\\&\qquad+Y\xi\xi-W\xi(1-x))
\end{aligned}
\]

Supponamus, ex divisione functionis $\xi X$ per

\[
\begin{aligned}
x^{p-1}+x^{p-2}+x^{p-3}+\text{etc.}+x+1
\end{aligned}
\]

oriri quotientem $U$ cum residuo $T$, sive haberi

\[
\begin{aligned}
\xi X=\frac{1-x^p}{1-x}\cdot U+T
\end{aligned}
\]

ita ut $U,T$ sint functiones integrae, etiam respectu coëfficientium numericorum, et quidem $T$ ordinis certe inferioris, quam divisor. Erit itaque

\[
\begin{aligned}
qT-\varepsilon p(\delta p^{\frac12(q-1)}-\gamma)\\&=\frac{1-x^p}{1-x}\cdot(Z(\delta p^{\frac12(q-1)}-\gamma)\\&\qquad+Y\xi\xi-W\xi(1-x)-qU)
\end{aligned}
\]

quae aequatio manifesto subsistere nequit, nisi tum membrum a laeva tum membrum a dextra per se evanescat. Erit itaque $\varepsilon p(\delta p^{\frac12(q-1)}-\gamma)$ per $q$ divisibi%
\SourceLeaf{69}%
lis, nec non etiam $\delta p^{\frac12(q-1)}-\gamma$, adeoque etiam propter $\delta\delta=1$, numerus $p^{\frac12(q-1)}-\gamma\delta$ per $q$ divisibilis erit.

Quodsi iam per $\beta$ designatur unitas positive vel negative accepta, prout $p$ est residuum vel non-residuum quadraticum numeri $q$, erit $p^{\frac12(q-1)}-\beta$ per $q$ divisibilis, adeoque etiam $\beta-\gamma\delta$, quod fieri nequit, nisi fuerit $\beta=\gamma\delta$. Hinc vero theorema fundamentale sponte sequitur. Scilicet

I. Quoties vel uterque $p,q$, vel alteruter tantum est formae $4k+1$, adeoque $\delta=+1$, erit $\beta=\gamma$, et proin vel simul $q$ residuum quadraticum ipsius $p$, atque $p$ residuum quadraticum ipsius $q$; vel simul $q$ non-residuum ipsius $p$, atque $p$ non-residuum ipsius $q$.

II. Quoties uterque $p,q$ est formae $4k+3$, adeoque $\delta=-1$, erit $\beta=-\gamma$, adeoque vel simul $q$ residuum quadraticum ipsius $p$, atque $p$ non-residuum ipsius $q$; vel simul $q$ non-residuum ipsius $p$, atque $p$ residuum ipsius $q$.

Q. E. D.

